% Lösungen Einheit 4 von 5 – Sachaufgaben (zusammenbau v0.1)
\einheitenkopf{Einheit 4 von 5 · Sachaufgaben}

%% TODO Lösung der Erklärzeile 23g) fehlt in der Bank
\erg{23}{a) rechts: $3{,}10\,€$ und $3{,}70\,€$; vor $x$ und $y$: $3$, $2$ und $1$, $4$; $x$ Preis eines Brötchens, $y$ Preis einer Brezel (in €) \quad b) I: $x + y = 4{,}60$, II: $3x + 4y = 15{,}80$ \quad c) I: $x + y = 21$, II: $2x + 5y = 69$ \quad d) I: $x + y = 3{,}40$, II: $4x + 2y = 8{,}80$ \quad e) I: $x + y = 45$, II: $2x + 3y = 110$ \quad f) I: $x + y = 14$, II: $2x + 5y = 46$ \quad g) \ldots \quad h) I: $x + y = 50$, II: $2{,}80x + 4{,}50y = 177{,}40$ \quad i) I: $x + y = 23$, II: $2x + 4y = 68$ \quad j) I: $x + y = 31$, II: $x - y = 7$ \quad k) $x$: Preis eines Blocks in €, $y$: Preis eines Füllers in € \quad l) Im Korb liegen zusammen $25$ Äpfel und Birnen. \quad m) I: $2x + y = 14{,}50$, II: $x + 2y = 12{,}50$; $x = 5{,}50$, $y = 3{,}50$. Ein Burger kostet $5{,}50\,€$, eine Portion Pommes $3{,}50\,€$. \quad n) I: $x + y = 4{,}10$, II: $5x + 2y = 13{,}60$}
%% TODO Lösung der Erklärzeile 24g) fehlt in der Bank
\erg{24}{a) eine Länge \quad b) I: Es werden zusammen $60$ Brote gebacken. II: Es sind doppelt so viele Roggenbrote wie Weizenbrote. \quad c) I: Es wurden zusammen $150$ Karten verkauft. II: Die Einnahmen betragen zusammen $1\,500\,€$. \quad d) I: Die Mischung wiegt zusammen $50$ kg. II: Es ist $30$ kg mehr Erde als Sand. \quad e) I: Der Umfang ist $36$ cm. II: Die eine Seite ist $4$ cm länger als die andere. \quad f) I: Die Mutter ist dreimal so alt wie Paul. II: Zusammen sind sie $48$ Jahre alt. \quad g) \ldots \quad h) II: Der Eiweißgehalt der Mischung ist $30\,\%$ – jede Sorte trägt ihren Gehalt mal ihren Anteil bei. I: Die Anteile ergeben zusammen die ganze Mischung. \quad i) I: Stickstoffbilanz – die Mischung hat $9\,\%$ Stickstoff, jeder Dünger trägt seinen Gehalt mal seinen Anteil bei. II: $u$, $v$, $w$ sind die Anteile der drei Dünger an der Mischung (keine Prozentzahlen); zusammen ergeben sie die ganze Mischung.}
\erg{25}{a) I: $x + y = 30$, II: $14x + 5y = 240$; $x = 10$, $y = 20$: $10$ kg Nüsse und $20$ kg Rosinen.}
\erg{26}{a) $5 + 4 = 9$ (wA); $3 \cdot 5 + 2 \cdot 4 = 23$ (wA): Die Lösung stimmt.}
\erg{27}{a) $0{,}09$: Preis einer Minute in €; $12{,}60$: Betrag der ganzen Rechnung in €.}
\erg{28}{a) Die Beträge sind vertauscht: richtig ist I: $x + y = 4{,}20$, II: $2x + 5y = 14{,}10$. \quad b) Zuordnung vertauscht: Zweibettzimmer $x = 8$, Dreibettzimmer $y = 9$. \quad c) Die Nussanteile in Prozent sind die Vorzahlen $5$, $15$ und $30$; $u$, $v$, $w$ sind die Anteile der drei Sorten an der Mischung – darum ist ihre Summe $1$. \quad d) Die Gleichung $x + y = 8$ hat viele Lösungspaare, etwa $3$ und $5$ oder $2$ und $6$; erst eine zweite Angabe wählt eines davon aus. \quad e) Die drei Anteile bilden zusammen die ganze Mischung, also $100\,\%$ oder $1$. Die zweite Gleichung zählt den Gehalt der Zutat: Was jede Sorte mit ihrem Anteil beiträgt, ergibt zusammen den Gehalt der Mischung.}
